\chapter{The odd--odd sector}
\label{chap:odd-odd}

Here $\psi$ vanishes on the separatrix, so the even--even boundary term is
replaced by a normal derivative term of the wrong sign. The fix is to divide
by $\Psi$. The quotient $\phi=\psi/\Psi$ is even across the separatrix and
satisfies a weighted problem. A weighted virial identity, a Picone
decomposition with an explicit comparison function, and a corrected
boundary form then give the contradiction.

\emph{Local notation (as in the manuscript):} on $\Cell$, $s=\Psi$,
$r=|\nabla\Psi|^2$ (not a distance), $\mu=s^2$, $\delta=\cos x-\cos y$,
$n=\nabla s/\sqrt r$, and $t$ the unit tangent.

\section{The quotient problem}

\begin{definition}[Odd--odd quotient]
\label{def:dd-quotient}
\uses{def:eigen-problem, def:cells}
For an odd--odd eigenfunction let $\phi=\psi/\Psi$ and $\nu=\omega/\Psi$, and
on $\Cell$ define $E=1-s^{-2}\Div(s^2\nabla)$.
\end{definition}

\begin{lemma}[Quotient equations]
\label{lem:dd-quotient-equations}
\uses{def:dd-quotient}
$\nu=-E\phi$. Multiplication by $s$ is unitary from $L^2(\Cell,\mu\,dx)$ onto
$L^2(\Cell)$ and conjugates $E$ to $-\Delta$: $sE(s^{-1}u)=-\Delta u$. If
$(\omega,\psi)$ is a transport pair then so is $(\nu,\phi)$, since $L_0\Psi=0$.
\end{lemma}

\begin{proof}
\uses{lem:dipole-identities}
$\Delta(s\phi)=s\Delta\phi+2\nabla s\cdot\nabla\phi-s\phi
=s\big(s^{-2}\Div(s^2\nabla\phi)-\phi\big)$, using $\Delta s=-s$. The
conjugation is the same computation.
\end{proof}

\begin{lemma}[Orthogonality and parity of the quotient]
\label{lem:dd-orthogonality}
\uses{def:dd-quotient}
$\int_\Cell\Psi^2\phi=0$, and $\phi$ is even under central inversion.
\end{lemma}

\begin{proof}
\uses{lem:H-selfadjoint, lem:dipole-identities, lem:stream-sector}
$H\Psi=0$, so pairing the eigenvalue equation with $\Psi$ gives
$\omega\perp\Psi$; then $\Delta\Psi=-\Psi$ gives $\psi\perp\Psi$. The product
$\psi\Psi$ is invariant under the reflections exchanging the cells, so
$\int_\Cell\Psi\psi=0$. Both $\psi$ and $\Psi$ are centrally even.
\end{proof}

\section{The weighted virial identity}

\begin{definition}[Odd--odd virial field]
\label{def:dd-virial-field}
\uses{def:dd-quotient}
Let $q(s)=(s^3-8)/(3s^2)$, $W=q(s)\nabla s/r$, and
\[
  F=\frac{s(8-s^3)\delta^2}{3r^2},\quad
  C_0=\frac{s(8-s^3)}{3r},\quad
  A_0=\frac{2s^2}3+\frac{32}{3s}-C_0,\quad
  M=A_0\,n\otimes n+C_0\,t\otimes t,
\]
and $P=-\Div(M\nabla)-F$. Note $\delta^2=4-2r-s^2$.
\end{definition}

\begin{lemma}[Odd--odd virial algebra]
\label{lem:dd-virial-algebra}
\uses{def:dd-virial-field, def:adapted-field}
$W$ is adapted with coefficient $f=F/\mu$, it has constant defect
$f-\Div_\mu W=-1$ where $\Div_\mu W=\mu^{-1}\Div(\mu W)$, and
$\mu(2\sym DW-fI)=M$.
\end{lemma}

\begin{proof}
\uses{lem:normal-field-algebra}
Write $W=\alpha(s)W_{\rm ee}$, where
$\alpha(s)=(s^2+2s+4)/(3s^2)$ and
$W_{\rm ee}=(s-2)\nabla s/r$. Since $L_0s=0$, the adapted coefficient is
$f=\alpha f_{\rm ee}=F/\mu$. Also
\[
 f-\Div_\mu W=-\alpha-(s-2)\alpha'-2q/s=-q'-2q/s=-1.
\]
In the normal--tangent frame, Lemma~\ref{lem:normal-field-algebra} gives
\[
 \mu(2\sym DW-fI)
 =\operatorname{diag}\bigl(2s^2q'-C_0,C_0\bigr),\qquad
 C_0=s^2\alpha\frac{s(2-s)}r.
\]
Now $q'=1/3+16/(3s^3)$, giving the asserted $A_0$.
\end{proof}

\begin{definition}[Weighted generator, bulk density and level term]
\label{def:dd-virial-densities}
\uses{def:dd-virial-field}
On the punctured positive cell put
\[
 L_\mu=\mu^{-1}\Div(\mu\nabla),\qquad
 d_\mu=\Div_\mu W,\qquad D_\mu=W\cdot\nabla+\tfrac12d_\mu.
\]
For a function $u$, wherever the expressions are defined, set
\[
 j[u]=F(|Eu|^2-|u|^2)+\ip{M\nabla u}{\nabla u},\qquad
 J_\Omega[u]=\int_\Omega j[u].
\]
With $p=\partial_nu$ and $n=\nabla s/\sqrt r$ pointing toward increasing $s$,
let the level term be
\[
 G_u(\sigma)=\int_{s=\sigma}\left[
 \frac{8-\sigma^3}{3\sqrt r}
 \bigl(|p|^2-|\partial_\tau u|^2-|u|^2+
             \operatorname{Re}(Eu\,\bar u)\bigr)
 -\sigma^2\operatorname{Re}(\bar u p)\right]d\ell.
\]
Here $d\ell=\sqrt r\,dt$ and $\tau$ is either unit tangent.
\end{definition}

\begin{lemma}[The weighted first identity on an annulus]
\label{lem:dd-weighted-first}
\uses{def:dd-virial-densities, def:eigen-problem}
For an odd--odd eigenfunction with $\lambda\ne0$ and real
$\chi\in C_c^\infty(0,2)$, put
$W_\chi=\chi(s)W$, $f_\chi=\chi(s)f$ and
$D_{\mu,\chi}=W_\chi\cdot\nabla+\tfrac12\Div_\mu W_\chi$. Then
\[
 \operatorname{Re}\int_\Cell\mu f_\chi\nu\overline{\nu+\phi}
 =2\operatorname{Re}\int_\Cell\mu D_{\mu,\chi}\phi\,\bar\nu.
\]
Only the original eigenfunction domain is required, not
Proposition~\ref{prop:dd-quotient-domain}.
\end{lemma}

\begin{proof}
\uses{lem:virial-l2, lem:dd-virial-algebra, lem:dd-quotient-equations}
Extend $W_\chi,f_\chi$ periodically and by zero outside the positive cell.
They are smooth, supported away from $s=0,2$, and adapted because $L_0s=0$.
For the ordinary half-density generator $D_\chi$ one has
\[
 \Div_\mu W_\chi=\Div W_\chi+2W_\chi\cdot\nabla\log s,
 \qquad sD_{\mu,\chi}\phi=D_\chi(s\phi)=D_\chi\psi.
\]
Substitute $\omega=s\nu$ and $\psi=s\phi$ in the first identity of
Lemma~\ref{lem:virial-l2}. All integrands have annular support, where division
by $s$ is smooth and bounded. No weighted mollifier or inverse of $E$ is
needed. The argument does not use reflection parity.
\end{proof}

\begin{lemma}[Local weighted Green identity]
\label{lem:dd-green}
\uses{def:dd-virial-densities}
For every complex $C^3$ function $u$ on the punctured cell,
\[
 \mu f\bigl(|Eu|^2-\operatorname{Re}(Eu\,\bar u)\bigr)
 +2\mu\operatorname{Re}(D_\mu u\,\overline{Eu})
 =j[u]+\Div\mathcal V[u],
\]
where the real vector field on the right is
\[
 \mathcal V[u]=\mu\left[-\operatorname{Re}(\bar u\nabla u)
 -2\operatorname{Re}\bigl((W\cdot\nabla u)\overline{\nabla u}\bigr)
 +W(|\nabla u|^2+|u|^2)\right].
\]
\end{lemma}

\begin{proof}
\uses{lem:dd-virial-algebra}
Expand $E=1-L_\mu$. The product rule gives
\[
 \mu f\operatorname{Re}(L_\mu u\,\bar u)
 =-\mu f|\nabla u|^2+\tfrac12\mu L_\mu f|u|^2
 +\Div\left(\mu f\operatorname{Re}(\bar u\nabla u)
                  -\tfrac12\mu\nabla f|u|^2\right)
\]
and
\begin{align*}
 -2\mu\operatorname{Re}(D_\mu u\,\overline{L_\mu u})
 ={}&2\mu\ip{\sym DW\nabla u}{\nabla u}
       -\tfrac12\mu L_\mu d_\mu|u|^2\\
 &+\Div\left[-2\mu\operatorname{Re}
                (D_\mu u\,\overline{\nabla u})
       +\mu W|\nabla u|^2+\tfrac12\mu\nabla d_\mu|u|^2\right].
\end{align*}
Finally $2\mu\operatorname{Re}(D_\mu u\,\bar u)=\Div(\mu W|u|^2)$.
The bulk zeroth-order term is
$\mu[\tfrac12L_\mu(f-d_\mu)-f]|u|^2=-F|u|^2$.
In the flux, $f-d_\mu=-1$ leaves exactly $\mathcal V[u]$.
\end{proof}

\begin{lemma}[Odd--odd annular virial]
\label{lem:dd-annular-virial}
\uses{def:dd-virial-densities, def:eigen-problem}
For an odd--odd eigenfunction with $\lambda\ne0$, $j[\phi]$ is integrable
on every compact level annulus. There is a null set $N\subset(0,2)$ such
that for all $0<a<b<2$ outside $N$,
\[
 J_{\{a<s<b\}}[\phi]=G_\phi(a)-G_\phi(b).
\]
\end{lemma}

\begin{proof}
\uses{lem:dd-weighted-first, lem:dd-green, lem:dd-quotient-equations,
 lem:inverse-laplacian-derivatives, lem:lift-sobolev-local,
 lem:coarea-level-cutoff, lem:du-bois-reymond}
On every compact annulus, $\phi=\psi/s$ has local $H^2$ regularity and
$E\phi=-\nu\in L^2$, by the original inverse-Laplacian regularity.
The Green identity holds distributionally there: its product rules involve
only products of $L^2$ functions and their first weak derivatives, with
at most two derivatives of $\phi$. In particular its flux and bulk terms
are locally integrable; no derivative of $\nu$ is used.

Since $D_{\mu,\chi}=\chi D_\mu+\tfrac12\chi'q$, substituting $\nu=-E\phi$
in Lemma~\ref{lem:dd-weighted-first} and then using
Lemma~\ref{lem:dd-green} gives
\[
 \int_\Cell\chi(s)j[\phi]
 =\int_\Cell\chi'(s)
    \left(\nabla s\cdot\mathcal V[\phi]
                   -\mu q\operatorname{Re}(E\phi\,\bar\phi)\right).
\]
The quantity in parentheses divided by $\sqrt r$ is precisely the
$d\ell$ integrand of $G_\phi(s)$: use
$W\cdot\nabla\phi=qp/\sqrt r$ and $\mu q=(s^3-8)/3$.
Coarea therefore gives
$\int\chi(s)m(s)\,ds=\int\chi'(s)G_\phi(s)\,ds$, where
$m(s)=\int_{\{\Psi=s\}}j[\phi]dt$; both functions are locally integrable.
The one-dimensional weak derivative identity is $G_\phi'=-m$.
Lemma~\ref{lem:du-bois-reymond} and coarea give the assertion outside a
single null set.
\end{proof}

\begin{lemma}[Odd--odd truncated virial]
\label{lem:dd-truncated-virial}
\uses{def:dd-virial-densities, def:eigen-problem}
For an odd--odd eigenfunction with $\lambda\ne0$ and almost every
$\epsilon\in(0,2)$, the density $j[\phi]$ is integrable on $\{s>\epsilon\}$
and
\[
 J_{\{s>\epsilon\}}[\phi]=G_\phi(\epsilon).
\]
\end{lemma}

\begin{proof}
\uses{lem:dd-annular-virial, lem:dd-virial-algebra, lem:normal-field-algebra,
 lem:normal-coeff-bound, lem:inverse-laplacian-derivatives,
 lem:continuous-representative, lem:orbit-coordinates-area, lem:period-integral}
On $\{s>\epsilon\}$ division by $s$ preserves $H^2$ regularity and $\phi$
is bounded. The representation $W=\alpha W_{\rm ee}$ in
Lemma~\ref{lem:dd-virial-algebra} shows that $F$ and $M$ are bounded there,
including as $s\to2$. Thus $j[\phi]$ is integrable.

Put $H(s)=\int_{\{\Psi=s\}}(|\phi|^2+|\nabla\phi|^2+|E\phi|^2)dt$ near
$s=2$. Coarea gives $H\in L^1$ near $2$. The period is bounded there,
$r\le2(2-s)$, and $0\le(8-s^3)/3\le4(2-s)$. In the flow-time expression
for $G_\phi$, Cauchy--Schwarz and the boundedness of $\phi$ give
\[
 |G_\phi(s)|\le C\left(\sqrt{(2-s)H(s)}+(2-s)H(s)\right)
\]
for almost every $s$ near $2$. Outside any prescribed null set there are
$b_j\to2$ with $(2-b_j)H(b_j)\to0$, since otherwise $H(s)$ would dominate
a positive multiple of $(2-s)^{-1}$. Hence $G_\phi(b_j)\to0$.
Apply Lemma~\ref{lem:dd-annular-virial} with upper level $b_j$ and use
dominated convergence. No smooth extension of $W$ through the center and
no higher eigenfunction regularity there are required.
\end{proof}

\begin{lemma}[Finite part of the bulk weight]
\label{lem:dd-finite-part}
\uses{def:dd-virial-field, def:period, lem:period-orbit, lem:period-elliptic}
With $T(s)=4\mathbf K(1-s^2/4)$ the period on the level $\{\Psi=s\}$,
\[
  \int_{\{s>\epsilon\}}F=\frac{8-\epsilon^3}3T(\epsilon)-\int_{\{s>\epsilon\}}s^2,
  \qquad
  \FP\int_\Cell F=\frac{32}3\log8-2\pi^2,
\]
where the finite part removes $\tfrac{32}{3}\log(1/\epsilon)$, that is
$\tfrac83$ per corner.
\end{lemma}

\begin{proof}
\uses{lem:orbit-coordinates-area, lem:period-regularity, lem:period-asymptotics}
Coarea and integration by parts; $\int_\Cell\Psi^2=2\pi^2$ and
$T(\epsilon)=4\log8+4\log(1/\epsilon)+o(1)$.
\end{proof}

\begin{proposition}[Renormalized virial identity]
\label{prop:dd-virial}
\uses{def:dd-virial-densities, lem:dd-quotient-equations}
\notready
For an odd--odd eigenfunction with $\lambda\ne0$, with the bulk form
$J_\Omega[\phi]=\int_\Omega j[\phi]$ and
\[
  J_{\FP}[\phi]=\lim_{\epsilon\to0}\Big\{J_{\{s>\epsilon\}}[\phi]
  +\tfrac83\sum_P|\phi(P)|^2\log(1/\epsilon)\Big\},
\]
one has $J_{\FP}[\phi]=-\tfrac83\sum_{\text{edges}}Q_\lambda$, where on each edge
$g=\phi|_{\text{edge}}=(P_t-\lambda)z$, $\nu=-P_tz$ and
\[
  Q_\lambda=\int\frac{|g'|^2}{2\sech^2t}-\operatorname{Re}\int P_tz\,\bar g
  +\int\big(|g|^2-|g_-|^2\mathbf1_{t<0}-|g_+|^2\mathbf1_{t>0}\big)
  +\frac{\log8}2(|g_-|^2+|g_+|^2).
\]
\end{proposition}

\begin{proof}
\uses{lem:dd-truncated-virial, lem:dd-finite-part, prop:dd-quotient-domain,
 lem:pairing-limit, lem:ee-gradient-traces}
Lemma~\ref{lem:dd-truncated-virial} supplies the exact boundary integral
$J_{C_\epsilon}[\phi]=G_\phi(\epsilon)$ for almost every $\epsilon$.
The remaining step is to pass to its renormalized limit: the proposed route
uses the even reflection of $\phi$ to eliminate the normal derivative,
canonical transport for $(\nu,\phi)$ to identify the pairing, and the
saddle values of $\phi$ to subtract the corner logarithms. The required
weighted convergence and finite-part matching remain to be justified.
\end{proof}

\begin{remark}[Digest note: weighted identity versus separatrix limit]
The manuscript's ``same finite-difference argument'' is unnecessary for
annular fields: Lemma~\ref{lem:dd-weighted-first} transfers the established
unweighted identity directly by multiplication by $s$.
Lemmas~\ref{lem:dd-green}--\ref{lem:dd-truncated-virial} derive its boundary
integral, including the center limit, without the quotient-domain theorem.
Here $n=\nabla s/\sqrt r$ is the inward normal of $C_\epsilon$, not its
outward normal; this fixes the sign of the term
$-\epsilon^2\operatorname{Re}(\bar\phi\partial_n\phi)$.

The unresolved part is the passage $G_\phi(\epsilon)\to-\tfrac83\sum Q_\lambda$
after renormalization. The quoted even--even gradient traces apply to a
different sector and do not by themselves prove the required weighted
quotient limits. In particular, even reflection gives zero normal trace
on a regular edge, but a uniform corner-tail estimate is also needed.
The quotient-domain argument itself remains subject to review.
\end{remark}

\section{The Picone decomposition}

\begin{definition}[Logarithmic ground state]
\label{def:dd-ground-state}
\uses{def:dd-virial-field}
Let
\[
  w_*=1+\frac{r+s^2}{16}-\frac{r^2}{256}\log(r/2),\qquad
  w_b(r)=1+\frac r{16}-\frac{r^2}{256}\log(r/2),
\]
(with $r^2\log r$ extended continuously at $0$), and
\[
  B(r)=1-\frac{4[r/8+(4-3r)w_b'(r)]}{w_b(r)}
  =\frac{5r/16+(r/256)[(32-25r)\log(r/2)+16-12r]}{w_b(r)}.
\]
\end{definition}

\begin{lemma}[Boundary current of the ground state]
\label{lem:dd-picone-current}
\uses{def:dd-ground-state, prop:dd-virial}
\notready
For an odd--odd eigenfunction,
\[
  \int_\Cell F|E\phi|^2+\int_\Cell Mw_*^2|\nabla(\phi/w_*)|^2
  +\int_\Cell\frac{Pw_*}{w_*}|\phi|^2
  =-\frac83\sum_{\text{edges}}\Big\{\mathcal F_\lambda[z]+\int B(r)|g|^2\,dt\Big\},
\]
with all bulk terms ordinary (absolutely convergent) integrals and
$\mathcal F_\lambda$ the first line of $Q_\lambda$ (manuscript
\texttt{(r4:picone-global)}).
\end{lemma}

\begin{proof}
\uses{lem:picone, prop:dd-virial, lem:dipole-identities, prop:dd-quotient-domain}
Apply Lemma~\ref{lem:picone} on $C_\epsilon$ with $g=w_*$. Using
$\nabla s\cdot\nabla r=s(4-s^2-3r)$, the boundary current in flow time is
$b_\epsilon|\phi|^2dt$ with $b_\epsilon=-\tfrac83(1-B(r))+O(\epsilon)$
uniformly. The universal $-\tfrac83$ part reproduces exactly the finite part
already present in $Q_\lambda$, so the corner logarithms cancel. Near a corner
$Pw_*=O(1+|\log r|)$, so the residual is integrable.
\end{proof}

\begin{lemma}[Bulk coefficient bounds]
\label{lem:dd-bulk-coefficients}
\uses{def:dd-ground-state}
On $\Cell$, writing $w=w_*$,
\[
  1\le w\le\tfrac54,\qquad Mw^2\ge\tfrac{25}{16}\mu I,\qquad Pw\ge-\tfrac52\mu w.
\]
\end{lemma}

\begin{proof}
\uses{lem:dipole-identities, cert:dd-residual}
The geometric range is $0<s<2$, $s(2-s)\le r\le2-s^2/2$. Since $0<r\le2$, the
logarithmic term is nonnegative, and $s^2+r\le4-r$, $r\log(2/r)\le2/e$ give
$w\le5/4$. Next, $3sr(A_0-C_0)\ge4s(s-2)^2(s+4)\ge0$, so it suffices to bound
$C_0w^2$; for fixed $s$, $(16+r+s^2)^2/r$ decreases in $r$, and at
$r=2-s^2/2$ the inequality reduces to
$2(s^2+2s+4)(s^2+36)^2-4800s(s+2)\ge0$, which with $t=2-s$ is
$t(1920+4544t-1984t^2+312t^3-28t^4+2t^5)\ge0$, true since
$4544-1984t\ge576$ and $312-28t\ge256$ on $[0,2]$. The bound
$Pw\ge-\tfrac52\mu w$ is Lemma~\ref{cert:dd-residual}.
\end{proof}

\begin{lemma}[Weighted Poincar\'e inequality]
\label{lem:zonal-poincare}
\uses{def:cells}
\notready
Normalize $\mu\,dx$ to a probability measure on $\Cell$. If $z$ is even under
central inversion, then
$\int\mu|\nabla z|^2\ge3\int\mu|z-\ip{z}{1}_\mu|^2$.
\end{lemma}

\begin{proof}
\uses{lem:cell-square, lem:separatrix-lines, lem:diagonal-calculus}
In the $(a,b)$ coordinates $\mu\,dx=c\cos^2a\cos^2b\,da\,db$, and the weighted
Neumann Laplacian is half the sum of two zonal Laplacians on $S^3$ (polar angle
$a+\pi/2$), with eigenvalues $n(n+2)$. Central inversion removes the modes of
degree $(1,0)$ and $(0,1)$; the next is $(1,1)$ with eigenvalue
$\tfrac12(3+3)=3$.
\end{proof}

\begin{remark}[Digest note]
Lemma~\ref{lem:zonal-poincare} needs a precise statement of the function space
(the form domain of the weighted Neumann Laplacian on the square) and of the
spectral decomposition, as the ``zonal $S^3$'' identification is only
sketched. A self-contained alternative: separation of variables in $a$ and
$b$ with Gegenbauer polynomials $C^{(1)}_n(\sin a)$.
\end{remark}

\begin{lemma}[Kantorovich-type bound]
\label{lem:kantorovich}
For a probability measure and a function $1\le w\le5/4$,
$\Var(w)/\ip{w}{1}^2\le1/80$.
\end{lemma}

\begin{proof}
Integrate $(w-1)(\tfrac54-w)\ge0$ to get $\Var(w)\le(m-1)(\tfrac54-m)$ with
$m=\ip{w}{1}$, and maximize $(m-1)(\tfrac54-m)/m^2$ (at $m=10/9$).
\end{proof}

\begin{lemma}[Bulk Picone bound]
\label{lem:dd-bulk-picone}
\uses{def:dd-ground-state}
If $\phi$ is centrally even with $\int_\Cell\mu\phi=0$, then
\[
  \int_\Cell Mw_*^2|\nabla(\phi/w_*)|^2+\int_\Cell\frac{Pw_*}{w_*}|\phi|^2
  \ge\frac{25}{54}\int_\Cell\mu|\phi|^2 .
\]
\end{lemma}

\begin{proof}
\uses{lem:dd-bulk-coefficients, lem:zonal-poincare, lem:kantorovich}
With $z=\phi/w_*$ and $m=\ip{z}{1}_\mu$, the constraint $\ip{w_*z}{1}_\mu=0$
and Cauchy--Schwarz give $|m|^2\le\tfrac1{80}\int\mu|z-m|^2$. Hence
$\int Mw_*^2|\nabla z|^2\ge\tfrac{25}{16}\cdot\tfrac{3}{1+1/80}\int\mu|z|^2
\ge\tfrac{80}{27}\int\mu|\phi|^2$ (using $|\phi|\le\tfrac54|z|$), and
$\tfrac{80}{27}-\tfrac52=\tfrac{25}{54}$.
\end{proof}

\section{The corrected boundary form}

\begin{definition}[Corrected boundary form]
\label{def:corrected-boundary-form}
\uses{def:boundary-form, def:dd-ground-state, def:h1-interval, lem:h1-interval-endpoints}
With $x=\tanh t$, $w=1-x^2$, let
\[
  \beta(x)=\frac{\tfrac34-\tfrac3{16}w+(\tfrac14-\tfrac{25}{64}w)\log w}
               {1+\tfrac18w-\tfrac1{64}w^2\log w}=\frac{B(2w)}w .
\]
For $f\in H^1(-1,1)$ with $wf''\in L^2$ and $f(1)=\kappa f(-1)$, put
$g=-iwf'-\lambda f$ and
$\mathcal B_\lambda(f)=\mathcal F_\lambda(f)+\int_{-1}^1\beta|g|^2\,dx$.
\end{definition}

\begin{lemma}[Pointwise lower bound for $\beta$]
\label{lem:beta-lower}
\uses{def:corrected-boundary-form}
$\beta\ge\tfrac12+\tfrac14\log w$ on $(-1,1)$ (manuscript
\texttt{(eq:beta-lower)}).
\end{lemma}

\begin{proof}
Multiplying by the positive denominator, the difference is
$\tfrac{1-w}4+\tfrac{w^2-54w}{128}\log w+\tfrac{w^2}{256}(\log w)^2\ge0$ for
$0<w\le1$.
\end{proof}

\begin{lemma}[A positive scalar form]
\label{lem:scalar-form-positive}
\uses{def:h1-interval}
For $g\in H^1(-1,1)$, $\int_{-1}^1\big(|g'|^2+(2+\log w)|g|^2\big)\,dx\ge0$,
with equality only for $g=0$.
\end{lemma}

\begin{proof}
\uses{lem:picone}
Use the Picone identity with $h=1+w^2\log w/8>0$, whose endpoint derivatives
vanish. Then
$-h''+(2+\log w)h=\tfrac12+\tfrac74w+(\tfrac32w+\tfrac14w^2)\log w
+\tfrac18w^2(\log w)^2\ge\tfrac12-\tfrac32e^{-13/6}-\tfrac1{8e}>\tfrac{17}{64}$.
\end{proof}

\begin{proposition}[Positivity of the corrected boundary form]
\label{prop:corrected-trace}
\uses{def:corrected-boundary-form}
For $\kappa\in\{1,-1\}$ and real $\lambda\ne0$, $\mathcal B_\lambda(f)>0$
whenever $f\ne0$ (manuscript \texttt{(prop:corrected-trace)}).
\end{proposition}

\begin{proof}
\uses{lem:beta-lower, lem:scalar-form-positive, lem:trace-sos-extension, lem:affine-splitting, lem:hardy-1d}
For $u\in H^1_0$ with $wu''\in L^2$, write $F_0=\mathcal F_\lambda(u)$ and
$g_u=-iwu'-\lambda u$; then $\|g_u'\|_2^2\le4F_0$ and $g_u\in H^1_0$.
If $\kappa=1$, write $f=a+u$; then
$\mathcal B_\lambda(f)\ge\tfrac14\int(|g'|^2+(2+\log w)|g|^2)$, which is positive
by Lemma~\ref{lem:scalar-form-positive}, and $g=0$ forces $f=0$ when
$\lambda\ne0$. If $\kappa=-1$, write $f=bx+u$; then
$\mathcal F_\lambda(f)=F_0+\lambda^2|b|^2$, and estimating the negative part of
$\beta$ near the endpoints (with $a=e^{-2}$) and the positive part on
$[-\tfrac12,\tfrac12]$ gives
$\mathcal B_\lambda(f)\ge\tfrac25F_0+(\tfrac1{10}+\tfrac45\lambda^2)|b|^2$.
\end{proof}

\begin{remark}[Digest note]
The $\kappa=-1$ estimate chains several explicit numerical constants
($203/1920$, $\lambda^2/96$, $16/\pi^2$, \dots) that the manuscript states
without intermediate steps. They should be re-derived; this is a good
candidate for a careful Lean proof, since it is one-dimensional.
\end{remark}

\section{Conclusion}

\begin{theorem}[Odd--odd exclusion]
\label{thm:dd-exclusion}
\uses{def:eigen-problem, def:sectors}
\lean{NoEmbeddedEigenvalues.odd_odd_exclusion}
\leanok
If $\omega$ solves the eigenvalue problem with $\lambda\ne0$ and
$S_1^*\omega=S_2^*\omega=-\omega$, then $\omega=0$ (manuscript
\texttt{(r4:dd-exclusion)}).
\end{theorem}

\begin{proof}
\uses{lem:rotation-characters, prop:dd-quotient-domain, lem:dd-orthogonality, lem:dd-picone-current, lem:dd-bulk-picone, prop:corrected-trace}
Fix an $R$-character $\kappa$. By Lemmas~\ref{lem:dd-picone-current}
and~\ref{lem:dd-bulk-picone},
\[
  \frac{25}{54}\int_\Cell\Psi^2|\phi|^2\le
  \int F|E\phi|^2+\int Mw_*^2|\nabla(\phi/w_*)|^2+\int\frac{Pw_*}{w_*}|\phi|^2
  =-\frac83\sum_{\text{edges}}\mathcal B_\lambda(f)\le0
\]
by Proposition~\ref{prop:corrected-trace}. Hence $\phi=0$, so $\psi=0$ and $\omega=0$.
\end{proof}
